\chapter{Implementation}
\label{ch:implementation}

\section{Introduction}
This chapter describes the implementation details of the Seasonal Job Matching Platform Mobile Application. It details the development environment and package configurations, provides a high-level overview of the folder hierarchy, and details the programming logic for the core application features. These include the authentication session interceptors, memory-optimized hybrid pagination scroll handlers, parallel service aggregation, and optimistic state synchronization controllers.

\section{Software Environment and Packages}
The mobile application is built using the Flutter SDK. Table~\ref{tab:development_packages} lists the primary open-source packages and libraries declared in the application's configuration (\texttt{pubspec.yaml}), along with their versions and roles.

\begin{table}[htbp]
\centering
\caption{Core Flutter Packages and Implementations}
\label{tab:development_packages}
\begin{tabular}{|l|l|p{7.5cm}|}
\hline
\textbf{Package Name} & \textbf{Version} & \textbf{Operational Purpose in Client} \\ \hline
\texttt{flutter\_riverpod} & \texttt{\^{}3.0.3} & Reactively manages UI states and coordinates dependency injection. \\ \hline
\texttt{dio} & \texttt{\^{}5.9.0} & HTTP client for network requests, supporting interceptors. \\ \hline
\texttt{flutter\_secure\_storage} & \texttt{\^{}9.2.2} & Encrypts and caches user sessions and authentication tokens. \\ \hline
\texttt{firebase\_messaging} & \texttt{\^{}15.2.5} & Connects the client to Firebase Cloud Messaging to receive push notifications. \\ \hline
\texttt{flutter\_local\_notifications} & \texttt{\^{}19.2.1} & Displays in-app notification banners when messages arrive in the background. \\ \hline
\texttt{freezed\_annotation} & \texttt{\^{}3.1.0} & Code generator for immutable data models and JSON serialization. \\ \hline
\texttt{google\_fonts} & \texttt{\^{}6.3.3} & Custom typography, utilizing the Inter and Outfit font packages. \\ \hline
\end{tabular}
\end{table}

\section{Folder Hierarchy}
The codebase is organized into key subdirectories under the main source directory (\texttt{lib/}). This layout supports the clean separation of UI components, state managers, and data repositories:
\begin{itemize}
    \item \texttt{lib/core}: Configures the base HTTP clients (Dio), authentication security interceptors, and navigation helper utilities.
    \item \texttt{lib/data}: Manages networking queries and implements the repository interfaces (e.g., \texttt{ApplicationsRepository}).
    \item \texttt{lib/domain}: Encapsulates standalone use case classes that represent business logic transactions.
    \item \texttt{lib/models}: Contains immutable data structures and data transfer objects (DTOs).
    \item \texttt{lib/providers}: Coordinates the Riverpod state notifiers, screen filters, and data sync controllers.
    \item \texttt{lib/screens}: Hosts the main page templates and layouts (e.g., \texttt{LoginScreen}, \texttt{LayoutScreen}).
    \item \texttt{lib/services}: Wraps raw endpoint configurations and provides low-level API communication.
    \item \texttt{lib/theme}: Configures color schemes, custom animations, and typographic decorators.
    \item \texttt{lib/widgets}: Reusable UI components (cards, badges, loading shimmers).
\end{itemize}

\section{Authentication and Session Interception}
Authentication sessions are secured using JWT tokens. A custom HTTP interceptor, \texttt{AuthInterceptor}, automatically appends the token from secure storage to the headers of outgoing requests.

If a session expires or is revoked, the server returns a \texttt{401 Unauthorized} or \texttt{403 Forbidden} response. This is caught by a global Dio error interceptor. Listing~\ref{lst:interceptor_logic} shows the implementation of this error handling and session validation flow.

\begin{lstlisting}[language=Java, caption={Dio Session Expiration Interceptor implementation}, label=lst:interceptor_logic]
dio.interceptors.add(
  InterceptorsWrapper(
    onError: (error, handler) async {
      final requestPath = error.requestOptions.path.toLowerCase();
      final method = error.requestOptions.method.toUpperCase();

      // Skip verification for public authentication endpoints
      final isAuth = requestPath.contains('auth/') ||
                     requestPath.contains('users/login') ||
                     requestPath.contains('users/signup') ||
                     (requestPath == 'users' && method == 'POST');

      if (isAuth || NavigationService().isOnAuthScreen()) {
        handler.next(error);
        return;
      }

      if (error.response?.statusCode == 401 || 
          error.response?.statusCode == 403) {
        
        // Prevent duplicate popup modals
        if (AuthDialogManager().isSessionExpiredHandled) {
          handler.next(error);
          return;
        }

        AuthDialogManager().markSessionExpiredHandled();
        await storage.clearToken();
        await storage.clearUserId();
        
        // Log out the user and trigger UI redirect
        ref.read(authProvider.notifier).logout(sessionExpired: true);
      }
      handler.next(error);
    },
  ),
);
\end{lstlisting}

\section{Hybrid Pagination and Memory Optimization}
To support large datasets without degradation in performance, the application implements a hybrid pagination strategy for job listings.
The UI tracks the scroll offset using a \texttt{NotificationListener} wrapping the list builder.
* **Automatic Fetching:** For the first 200 items (4 pages of 50 listings each), the system automatically loads the next page when the user scrolls near the bottom of the list.
* **Manual Fetching:** Beyond 200 items, automatic scrolling pauses, and the user must click a manual ``Load More'' button to fetch additional pages.

This approach balances network overhead and memory usage. While a textual model occupies only $\approx 1.5\text{ KB}$ in memory, inflating that data into layout components (widgets, images, and haptics) within Flutter's render tree increases the resource cost. Restricting automatic scrolling to 200 items avoids memory issues on low-end devices.

\section{Parallel Aggregation and Optimistic Updates}
\subsection{Parallel Aggregation}
When initializing the applicant's profile screen, the system fetches data concurrently from four distinct endpoints to reduce wait times. This parallel request flow is managed within \texttt{PersonalInformationAsyncNotifier} using Dart's asynchronous libraries, as shown in Listing~\ref{lst:parallel_fetching}.

\begin{lstlisting}[language=Java, caption={Parallel request aggregation flow}, label=lst:parallel_fetching]
Future<PersonalInformationModel> fetchProfileData(String userId) async {
  final results = await Future.wait([
    _service.fetchUserData(userId),
    _service.fetchAppliedJobIds(userId),
    _service.fetchFavoriteJobIds(userId),
    _service.fetchFieldsOfInterest(userId),
  ]);

  return PersonalInformationModel(
    user: results[0] as UserModel,
    appliedJobIds: results[1] as List<int>,
    favoriteJobIds: results[2] as List<int>,
    fieldsOfInterest: results[3] as List<String>,
  );
}
\end{lstlisting}

\subsection{Optimistic Updates and Rollback}
To ensure the UI remains responsive, user actions such as favoriting a job are handled using optimistic state updates.
When a user toggles a favorite, the \texttt{FavoritesController} updates the local list in memory immediately, and the heart icon changes state before the network request completes.
If the API call fails, the controller catches the error, rolls back the local state to its previous version, and alerts the user. Listing~\ref{lst:optimistic_updates} shows the structure of this rollback loop.

\begin{lstlisting}[language=Java, caption={Optimistic updates and rollback loop implementation}, label=lst:optimistic_updates]
Future<void> toggleFavorite(int jobId) async {
  final previousState = state;
  
  // Optimistically update the UI state immediately
  state = _calculateToggledState(jobId);

  try {
    await _service.updateFavoriteJobs(userId, state.favoriteJobIds);
    // Refresh the detailed listings cache on success
    ref.invalidate(favoriteJobsProvider);
  } catch (e) {
    // Roll back to the original state on failure
    state = previousState;
    AppLogger.error('Failed to update favorites: $e');
    rethrow;
  }
}
\end{lstlisting}
